\documentclass[11pt,a4paper]{article}
\usepackage[T1]{fontenc}
\usepackage{lmodern,amsmath,amssymb,booktabs,geometry,hyperref,microtype}
\geometry{margin=27mm}
\title{When Does a More Complex Change-Point Detector Justify Its Cost?\\A Preregistered Monte Carlo Comparison of CUSUM, GLR, and BOCPD}
\author{Diogo Ribeiro}
\date{26 October 2026}
\begin{document}
\maketitle
\begin{abstract}
An alarm does not determine whether a model should be retrained. We compare the economic consequences of retraining policies driven by classical CUSUM, a Gaussian generalised likelihood-ratio (GLR) detector, and Bayesian online change-point detection (BOCPD). In the preregistered synthetic experiment, 199 independent null-calibration streams established detector thresholds; 100 evaluation streams per regime, each with 120 deployment observations, were used to assess the policies. Under gradual concept drift, mean economic regret was 0.217649 for CUSUM, 0.217060 for GLR and 0.226625 for BOCPD. The GLR--CUSUM regret difference had an adjusted paired interval of $[-0.004824,0.003344]$, which does not establish the specified minimum improvement of 0.01 regret units. BOCPD had higher regret (difference 0.008977; adjusted interval $[0.003687,0.013818]$). Neither candidate verified the required 10\% false-alarm ceiling. The data do not justify replacing CUSUM under this monitoring score, drift mechanism and cost specification.
\end{abstract}
\section{Introduction}
Sequential detection is often assessed through alarm rates and detection delay. Neither quantity determines whether acting on an alarm improves an economic decision. That depends on the retraining policy, its cost and the actions taken from the resulting predictions.

We compare classical cumulative sum (CUSUM) \cite{page1954}, Gaussian generalised likelihood ratios \cite{basseville1993}, and Bayesian online change-point detection (BOCPD) \cite{adams2007}. Each detector receives the same scalar monitoring signal. We evaluate their associated model-maintenance policies against prespecified economic-regret and false-alarm criteria; no new detection algorithm is proposed.
\section{Predictive and economic model}
Let $X_t\in\mathbb R^{16}$ and write $m_t(x)=\mathbb E[Y_t\mid X_t=x]$ for the synthetic conditional mean, with $\widehat m_t$ denoting the deployed regression estimate. The policy takes an action at cost $c$ when $\widehat m_t(X_t)>c$; refitting at time $t$ incurs cost $q$. Relative to an oracle with access to the complete conditional mean, one-step regret is
\begin{equation}
R_t=(m_t(X_t)-c)_+-(m_t(X_t)-c)\mathbf1\{\widehat m_t(X_t)>c\}+q\mathbf1\{\text{refit at }t\}.
\end{equation}
Regret accounts for action errors relative to the oracle and for model-maintenance costs. We use the true conditional mean only to evaluate simulated decisions. Neither the regression fit nor the detector observes it; the specification does not identify a causal treatment effect.

Eight latent factors generate the 16 predictors as correlated measurement pairs. The primary low-variance signal is proportional to $X_{t1}-X_{t2}$, whose coefficient gradually changes sign during deployment. This design isolates a specified change in the response relation. It was not fitted to agricultural, financial or healthcare observations.
\section{Sequential detectors}
We construct $Z_t$ from standardised disagreement between newly revealed outcomes and predictions of a frozen reference regression. CUSUM, GLR and BOCPD receive this same scalar sequence. The comparison therefore concerns three procedures applied to one monitoring statistic, not general multivariate detection.
\subsection{CUSUM}
For CUSUM, we use the one-sided recursion
\begin{equation}
G_t=\max\{0,G_{t-1}+Z_t-\kappa\}.
\end{equation}
The recursion accumulates evidence of positive departures in $Z_t$. Classical optimality results require a specified observation model and change hypothesis; they do not transfer automatically to a regression-derived score \cite{page1954,lorden1971}.
\subsection{Generalised likelihood ratio}
For the one-sided Gaussian GLR, we scan recent suffixes with lengths $m\in[m_{\min},m_{\max}]$:
\begin{equation}
\mathrm{GLR}_t=\max_m\frac{\left(\max\{0,\sum_{j=t-m+1}^t Z_j\}\right)^2}{2m}.
\end{equation}
This is the maximised Gaussian log likelihood ratio for an unknown positive mean shift at known unit variance. Serial dependence or a misspecified scale changes its reference distribution and may change detection performance.
\subsection{Bayesian online change-point detection}
BOCPD updates the posterior distribution of the run length $r_t$ under a constant change hazard, a conjugate Gaussian prior for the segment mean and known observation variance \cite{adams2007}. We use the posterior mass assigned to short run lengths as the alarm statistic. Its interpretation depends on that likelihood, hazard and prior; it is not a model-free probability of physical change.
\section{Preregistered comparison}
For each detector, we computed its maximum score over the monitoring horizon on 199 independent no-change calibration streams, then applied a corrected order-statistic threshold. The corresponding marginal false-alarm statement requires exchangeability of complete null-stream maxima over the same horizon. It does not apply to monitoring over an unrestricted time horizon.

The preregistered primary case fixed the low-variance response signal, a gradual-drift evaluation regime, immediate feedback, action cost $c=0.6$, refit cost $q=0.05$ and 120 deployment steps. We generated 100 independent evaluation streams per regime and compared GLR and BOCPD separately with CUSUM. A mean regret reduction of at least 0.01 per deployment observation was the minimum economically meaningful gain. The false-alarm ceiling was 10\%; the nominal 5\% family error budget was divided between paired economic comparisons and false-alarm bounds, with adjustment for the two candidates.

For candidate $A$, define the stream-paired effect
\[
\Delta_r(A)=\overline R_r(A)-\overline R_r(\mathrm{CUSUM}).
\]
A negative $\Delta_r(A)$ favours the candidate. Bootstrap limits are based on resampling independent streams, not adjacent time points. They describe Monte Carlo uncertainty conditional on the realised detection thresholds and exclude the variability of repeating the null calibration.
\section{Results}
\begin{table}[ht]
\centering
\caption{Mean net economic regret per observation under gradual drift.}
\begin{tabular}{lrr}
\toprule
Maintenance policy & Mean regret & Difference vs CUSUM\\
\midrule
Frozen & 0.249472 & +0.031824\\
Scheduled & 0.218668 & +0.001019\\
CUSUM & 0.217649 & 0\\
GLR & 0.217060 & $-0.000589$\\
BOCPD & 0.226625 & +0.008977\\
\bottomrule
\end{tabular}
\end{table}
The mean paired regret difference for GLR relative to CUSUM was $-0.000589$; its adjusted bootstrap limits were $[-0.004824,0.003344]$. The interval includes zero and does not establish the preregistered 0.01 improvement. For BOCPD, the difference was $+0.008977$, with adjusted limits $[0.003687,0.013818]$. Both limits are positive, indicating higher economic regret under this simulation.

\begin{table}[ht]
\centering
\caption{Probability of at least one alarm on 100 stable no-change streams.}
\begin{tabular}{lrr}
\toprule
Detector & Empirical alarm rate & Simultaneous upper bound\\
\midrule
CUSUM & 11\% & Not part of candidate eligibility\\
GLR & 6\% & 13.71\%\\
BOCPD & 10\% & 18.75\%\\
\bottomrule
\end{tabular}
\end{table}
The simultaneous upper false-alarm bounds for GLR and BOCPD exceeded the stipulated 10\% ceiling. The bounds do not show that their true false-alarm probabilities exceed 10\%; they leave compliance unverified. Both candidates received the preregistered verdict \texttt{material\_gain\_not\_supported}.
\section{Discussion and limitations}
The comparison does not determine a universally preferred detector. GLR's numerical advantage was smaller than the stipulated minimum material gain and was unresolved relative to zero. Scheduled retraining produced regret close to that of CUSUM. These estimates do not isolate which alarms or refits account for the remaining policy differences.

All detectors operate on one residual-derived score. We cannot infer how their relative performance would change with another monitoring statistic or with a joint model of the predictor distribution. The experiment also fixes one response mechanism, intervention cost, refit cost and horizon. Because the coefficient reverses gradually, there is no unique instantaneous change point against which to judge detection delay.

Detector parameters were not optimised for all possible shifts. Selecting a new configuration after inspecting these evaluation streams would require an independent confirmatory experiment. The bootstrap limits condition on the calibrated thresholds, and the null-stream bounds require the specified exchangeability assumption. The experiment uses a synthetic payoff, not an identified causal benefit.

A different run-length likelihood or an alternative detection score could change the ordering. The present experiment therefore supports a decision-level comparison under stated assumptions, not an ordering of algorithms independent of their application.
\section{Conclusion}
Neither GLR nor BOCPD satisfied the preregistered economic-gain and false-alarm evidence requirements for replacing CUSUM in the tested gradual-drift scenario. GLR's observed reduction in regret was too small and imprecisely estimated; BOCPD increased regret. The result concerns the specified monitoring score, model-maintenance policy and costs, and does not establish general optimality of CUSUM.
\section*{Reproducibility and availability}
The simulation code, preregistered protocol and checked numerical summary are available at
\href{https://github.com/DiogoRibeiro7/Medium-Blog/tree/main/statistical-evidence-and-decisions}{Medium-Blog/statistical-evidence-and-decisions}.
The full evaluation was executed as
\href{https://github.com/DiogoRibeiro7/Medium-Blog/actions/runs/40089779139}{GitHub Actions run 40089779139}.
The archived report is in \texttt{results/2026-10-24-detector-evidence/}. The complete tables are available as a time-limited workflow artifact, with their SHA-256 checksums recorded for verification.
\begin{thebibliography}{9}
\bibitem{page1954} Page, E. S. (1954). Continuous Inspection Schemes. \emph{Biometrika}, 41(1/2), 100--115.
\bibitem{basseville1993} Basseville, M., and Nikiforov, I. V. (1993). \emph{Detection of Abrupt Changes: Theory and Application}. Prentice Hall.
\bibitem{adams2007} Adams, R. P., and MacKay, D. J. C. (2007). Bayesian Online Changepoint Detection. arXiv:0710.3742.
\bibitem{lorden1971} Lorden, G. (1971). Procedures for Reacting to a Change in Distribution. \emph{The Annals of Mathematical Statistics}, 42(6), 1897--1908.
\end{thebibliography}
\end{document}
